\documentclass[10pt,a4paper]{amsart}
\usepackage{amsmath,amssymb,amsthm,graphicx}
\usepackage{verbatim}
\newenvironment{reference}{\comment}{\endcomment}
%\newenvironment{reference}{}{}
\newenvironment{slogan}{\comment}{\endcomment}
\newenvironment{history}{\comment}{\endcomment}

% For commutative diagrams we use Xy-pic
\usepackage[all]{xy}

% We use 2cell for 2-commutative diagrams.
\xyoption{2cell}
% Theorem environments.
%
\theoremstyle{plain}
\newtheorem{theorem}[subsection]{Theorem}
\newtheorem{proposition}[subsection]{Proposition}
\newtheorem{lemma}[subsection]{Lemma}

\theoremstyle{definition}
\newtheorem{definition}[subsection]{Definition}
\newtheorem{example}[subsection]{Example}
\newtheorem{exercise}[subsection]{Exercise}
\newtheorem{situation}[subsection]{Situation}

\theoremstyle{remark}
\newtheorem{remark}[subsection]{Remark}
\newtheorem{remarks}[subsection]{Remarks}

\numberwithin{equation}{subsection}

% Macros
%
\def\lim{\mathop{\mathrm{lim}}\nolimits}
\def\colim{\mathop{\mathrm{colim}}\nolimits}
\def\Spec{\mathop{\mathrm{Spec}}}
\def\Hom{\mathop{\mathrm{Hom}}\nolimits}
\def\Ext{\mathop{\mathrm{Ext}}\nolimits}
\def\SheafHom{\mathop{\mathcal{H}\!\mathit{om}}\nolimits}
\def\SheafExt{\mathop{\mathcal{E}\!\mathit{xt}}\nolimits}
\def\Sch{\mathit{Sch}}
\def\Mor{\mathop{\mathrm{Mor}}\nolimits}
\def\Ob{\mathop{\mathrm{Ob}}\nolimits}
\def\Sh{\mathop{\mathit{Sh}}\nolimits}
\def\NL{\mathop{N\!L}\nolimits}
\def\CH{\mathop{\mathrm{CH}}\nolimits}
\def\proetale{{pro\text{-}\acute{e}tale}}
\def\etale{{\acute{e}tale}}
\def\QCoh{\mathit{QCoh}}
\def\Ker{\mathop{\mathrm{Ker}}}
\def\Im{\mathop{\mathrm{Im}}}
\def\Coker{\mathop{\mathrm{Coker}}}
\def\Coim{\mathop{\mathrm{Coim}}}

% Boxtimes
%
\DeclareMathSymbol{\boxtimes}{\mathbin}{AMSa}{"02}

%
% Macros for moduli stacks/spaces
%
\def\QCohstack{\mathcal{QC}\!\mathit{oh}}
\def\Cohstack{\mathcal{C}\!\mathit{oh}}
\def\Spacesstack{\mathcal{S}\!\mathit{paces}}
\def\Quotfunctor{\mathrm{Quot}}
\def\Hilbfunctor{\mathrm{Hilb}}
\def\Curvesstack{\mathcal{C}\!\mathit{urves}}
\def\Polarizedstack{\mathcal{P}\!\mathit{olarized}}
\def\Complexesstack{\mathcal{C}\!\mathit{omplexes}}
% \Pic is the operator that assigns to X its picard group, usage \Pic(X)
% \Picardstack_{X/B} denotes the Picard stack of X over B
% \Picardfunctor_{X/B} denotes the Picard functor of X over B
\def\Pic{\mathop{\mathrm{Pic}}\nolimits}
\def\Picardstack{\mathcal{P}\!\mathit{ic}}
\def\Picardfunctor{\mathrm{Pic}}
\def\Deformationcategory{\mathcal{D}\!\mathit{ef}}
\usepackage{cleveref}
\crefname{originalreference}{original reference}{original references}
\Crefname{originalreference}{Original reference}{Original references}
\makeatletter
\newlabel{lemma-cap-c1-bivariant}{{0B77}{source}{}{sourcepointer-1}{}}
\newlabel{lemma-cap-c1-bivariant@cref}{{[originalreference][0][]0B77}{[1][0][]source}{}{}{}}
\makeatother
\crefname{originalreference}{original reference}{original references}
\Crefname{originalreference}{Original reference}{Original references}
\makeatletter
\newlabel{lemma-cap-projective-bundle}{{02TW}{source}{}{sourcepointer-2}{}}
\newlabel{lemma-cap-projective-bundle@cref}{{[originalreference][0][]02TW}{[1][0][]source}{}{}{}}
\makeatother
\crefname{originalreference}{original reference}{original references}
\Crefname{originalreference}{Original reference}{Original references}
\makeatletter
\newlabel{lemma-chow-ring-projective-bundle}{{02TX}{source}{}{sourcepointer-3}{}}
\newlabel{lemma-chow-ring-projective-bundle@cref}{{[originalreference][0][]02TX}{[1][0][]source}{}{}{}}
\makeatother
\crefname{originalreference}{original reference}{original references}
\Crefname{originalreference}{Original reference}{Original references}
\makeatletter
\newlabel{lemma-flat-pullback-bivariant}{{0B78}{source}{}{sourcepointer-4}{}}
\newlabel{lemma-flat-pullback-bivariant@cref}{{[originalreference][0][]0B78}{[1][0][]source}{}{}{}}
\makeatother
\crefname{originalreference}{original reference}{original references}
\Crefname{originalreference}{Original reference}{Original references}
\makeatletter
\newlabel{lemma-push-proper-bivariant}{{0EPK}{source}{}{sourcepointer-5}{}}
\newlabel{lemma-push-proper-bivariant@cref}{{[originalreference][0][]0EPK}{[1][0][]source}{}{}{}}
\makeatother
\crefname{originalreference}{original reference}{original references}
\Crefname{originalreference}{Original reference}{Original references}
\makeatletter
\newlabel{situation-setup}{{02QL}{source}{}{sourcepointer-6}{}}
\newlabel{situation-setup@cref}{{[originalreference][0][]02QL}{[1][0][]source}{}{}{}}
\makeatother
\title{Chern classes of a finite locally free sheaf are the unique bivariant classes satisfying the projective bundle relation}
\author{The Stacks Project Authors}
\date{stacks-project:spaces-chow.tex:local-snapshot-20261012; GNU FDL 1.2 or later}
\begin{document}
\maketitle
\noindent\textit{Editorial note.} All delivered mathematical text is the author's own, copied from the pinned Stacks Project chapter file spaces-chow.tex; no statement or proof step was rewritten, reordered or removed.\par\medskip
\noindent\textit{Editorial note.} The counted proof is the author's complete argument: a uniqueness argument from the projective bundle formula followed by an inductive construction of the Chern classes out of the bivariant Segre classes\par\medskip
\noindent\textit{Editorial note.} The label lemma-segre-classes names the counted statement, and the proof environment immediately following it in the source is delivered verbatim as its proof.\par\medskip
\noindent\textit{Editorial note.} The counted statement is carried by the author's own lemma environment in the chapter source.\par\medskip
\noindent\textit{Allowed standard prerequisites (editorial).}\begin{itemize}
\item The projective bundle formula, describing the Chow groups of a projective space bundle as a free module over those of the base, and the computation of the pushforward of powers of the first Chern class of the tautological quotient
\item The notion of a bivariant class, and the closure of bivariance under capping with a first Chern class, under flat pullback and under proper pushforward
\item The Segre class construction assigning to every base change of the projective bundle a map by pushing forward a power of the Chern class capped with the pullback
\end{itemize}\medskip
\ifdefined\hypertarget\hypertarget{sourcepointer-1}{}\fi
\noindent\textit{Source reference.} Original reference 0B77: sources/stacks-spaces-chow/pinned-original-source.tex, source lines 4144--4144. Prerequisite pointer invoked at proof line(s) 4633; it names a standard fact of the same chapter that the counted argument uses, and every step of the counted proof built on it is retained verbatim in the delivered body.\par\medskip
\ifdefined\hypertarget\hypertarget{sourcepointer-2}{}\fi
\noindent\textit{Source reference.} Original reference 02TW: sources/stacks-spaces-chow/pinned-original-source.tex, source lines 4321--4321. Prerequisite pointer invoked at proof line(s) 4636; it names a standard fact of the same chapter that the counted argument uses, and every step of the counted proof built on it is retained verbatim in the delivered body.\par\medskip
\ifdefined\hypertarget\hypertarget{sourcepointer-3}{}\fi
\noindent\textit{Source reference.} Original reference 02TX: sources/stacks-spaces-chow/pinned-original-source.tex, source lines 4381--4381. Prerequisite pointer invoked at proof line(s) 4622; it names a standard fact of the same chapter that the counted argument uses, and every step of the counted proof built on it is retained verbatim in the delivered body.\par\medskip
\ifdefined\hypertarget\hypertarget{sourcepointer-4}{}\fi
\noindent\textit{Source reference.} Original reference 0B78: sources/stacks-spaces-chow/pinned-original-source.tex, source lines 4160--4160. Prerequisite pointer invoked at proof line(s) 4634; it names a standard fact of the same chapter that the counted argument uses, and every step of the counted proof built on it is retained verbatim in the delivered body.\par\medskip
\ifdefined\hypertarget\hypertarget{sourcepointer-5}{}\fi
\noindent\textit{Source reference.} Original reference 0EPK: sources/stacks-spaces-chow/pinned-original-source.tex, source lines 4194--4194. Prerequisite pointer invoked at proof line(s) 4635; it names a standard fact of the same chapter that the counted argument uses, and every step of the counted proof built on it is retained verbatim in the delivered body.\par\medskip
\ifdefined\hypertarget\hypertarget{sourcepointer-6}{}\fi
\noindent\textit{Source reference.} Original reference 02QL: sources/stacks-spaces-chow/pinned-original-source.tex, source lines 264--264. Prerequisite pointer invoked at proof line(s) 4598; it names a standard fact of the same chapter that the counted argument uses, and every step of the counted proof built on it is retained verbatim in the delivered body.\par\medskip
\begin{lemma}
\label{lemma-segre-classes}
In Situation \ref{situation-setup} let $X/B$ be good.
Let $\mathcal{E}$ be a finite locally free sheaf of rank $r$ on $X$.
Let $(\pi : P \to X, \mathcal{O}_P(1))$ be the projective space
bundle associated to $\mathcal{E}$. For every
morphism $X' \to X$ of good algebraic spaces over $B$
there are unique maps
$$
c_i(\mathcal{E}) \cap - : \CH_k(X') \longrightarrow \CH_{k - i}(X'),\quad
i = 0, \ldots, r
$$
such that for $\alpha \in \CH_k(X')$ we have
$c_0(\mathcal{E}) \cap \alpha = \alpha$ and
$$
\sum\nolimits_{i = 0, \ldots, r}
(-1)^i c_1(\mathcal{O}_{P'}(1))^i \cap
(\pi')^*\left(c_{r - i}(\mathcal{E}) \cap \alpha\right) = 0
$$
where $\pi' : P' \to X'$ is the base change of $\pi$.
Moreover, these maps define a bivariant class
$c_i(\mathcal{E})$ of degree $i$ on $X$.
\end{lemma}

\begin{proof}
Uniqueness and existence of the maps $c_i(\mathcal{E}) \cap -$
follows immediately from Lemma \ref{lemma-chow-ring-projective-bundle}
and the given description of $c_0(\mathcal{E})$. For every $i \in \mathbf{Z}$
the rule which to every morphism $X' \to X$ of good algebraic spaces
over $B$ assigns the map
$$
t_i(\mathcal{E}) \cap - :
\CH_k(X') \longrightarrow \CH_{k - i}(X'),\quad
\alpha \longmapsto
\pi'_*(c_1(\mathcal{O}_{P'}(1))^{r - 1 + i} \cap (\pi')^*\alpha)
$$
is a bivariant class\footnote{Up to signs these are
the Segre classes of $\mathcal{E}$.} by Lemmas \ref{lemma-cap-c1-bivariant},
\ref{lemma-flat-pullback-bivariant}, and
\ref{lemma-push-proper-bivariant}.
By Lemma \ref{lemma-cap-projective-bundle} we have
$t_i(\mathcal{E}) = 0$ for $i < 0$ and $t_0(\mathcal{E}) = 1$.
Applying pushforward to the equation in the statement of the lemma
we find from Lemma \ref{lemma-cap-projective-bundle} that
$$
(-1)^r t_1(\mathcal{E}) + (-1)^{r - 1}c_1(\mathcal{E}) = 0
$$
In particular we find that $c_1(\mathcal{E})$ is a bivariant class.
If we multiply the equation in the statement of the lemma by
$c_1(\mathcal{O}_{P'}(1))$ and push the result forward to $X'$
we find
$$
(-1)^r t_2(\mathcal{E}) +
(-1)^{r - 1} t_1(\mathcal{E}) \cap c_1(\mathcal{E}) +
(-1)^{r - 2} c_2(\mathcal{E}) = 0
$$
As before we conclude that $c_2(\mathcal{E})$ is a bivariant class.
And so on.
\end{proof}

\end{document}
